\chapter{Materials}

\section{Linear Elasticity}

\section{Finite Elasticity}

\subsection{Kirchhoff-St.~Venant}

\subsection{Rivlin}

Let $I_1, I_2, I_3$ be the three invariants of the two Cauchy-Green 
deformation tensors.  The invariants are defined in terms of principal
stretch ratios $\lambda_1, \lambda_2, \lambda_3$,
\begin{align}
I_1 & = \lambda_1^2 + \lambda_2^2 + \lambda_3^2, \\
I_2 & = \lambda_1^2 \lambda_2^2 + \lambda_2^2 \lambda_3^2 + \lambda_3^2 \lambda_1^2, \\
I_3 & = \lambda_1^2 \lambda_2^2 \lambda_3^2.
\end{align}
Rivlin (1956) specified a strain energy density function $W$ as 
power series of the invariants,
\be
 W = \sum_{i+j+k=1}^{\infty} C_{ijk} (I_1 - 3)^i \; (I_2 - 3)^j \; (I_3 - 1)^k.
\ee
\begin{Hremark}
Note the presence of the $-3, -3$, and $-1$ terms in parenthesis allow for
zero strain energy density when there is zero stretch.
\end{Hremark}

\subsubsection{Incompressible Rivlin}

Next, {\bf assume incompressibility} of the material response.  In this case, 
$I_3 = 1$, and the strain energy density becomes
\be
 W = C_{10} (I_1 - 3) + C_{01} (I_2 - 3),  
\ee
which is called the {\bf Mooney-Rivlin equation}.

\subsubsection{Incompressible Neo-Hookean}

Using only the first term of the Mooney-Rivlin equation produces the 
{\bf Incompressible Neo-Hookean} material,
\be
 W = C_{10} (I_1 - 3).
\ee
Now, the invariants in terms of the right Cauchy-Green tensor $\tCbold$ are
\begin{align}
	I_1 & = \trace(\tCbold)  = C_{KK}, \\
	I_2 & = \half \left(I_1^2 - \tCbold \ddiv \tCbold \right) = 
	        \half \left(I_1^2 - C_{IK} C_{KI} \right), \\
	I_3 & = \det(\tCbold) = \det C_{IJ}.
\end{align}
Then, for small strains and Mooney-Rivlin,
\begin{align}
	E & = 6 \left(C_{10} + C_{01} \right), \\
	G & = 2 \left(C_{10} + C_{01} \right).
\end{align}
For small strains and Neo-Hookean 
\begin{align}
	C_{01} & = 0, \\
	C_{10} & = G/2, \implies E = 3 G. 
\end{align}

\subsubsection{Compressible Rivlin}
Next, allow the material to be compressible, characterized through 
bulk modulus $K$.  Thus, $I_3 \neq 1$.
The {\bf Mooney-Rivlin equation} for compressible materials is given by
\be
 W = W_{\smDEV}(I_3) +  C_{10} (I_1 - 3) + C_{01} (I_2 - 3),  
\ee
where $W_{\smDEV}$ is the hydrostatic work term.  

To fully uncouple the volumetric from the deviatoric responses, let the
{\bf volume preserving (isochoric) distortion invariants} be calculated
as
\begin{align}
	\bar{I}_1 & = I_1 \; I_3^{(-1/3)}, \\
	\bar{I}_2 & = I_2 \; I_3^{(-2/3)}, \\
	J & = \sqrt{I_3}.
\end{align}
Then, 
\be
 W = \underset{\smVOL}{\underbrace{\frac{K}{2}(J - 1)^2}} + 
 \underset{\smDEV}{\underbrace{C_{10} (\bar{I}_1 - 3) + C_{01} (\bar{I}_2 - 3)}},  
\ee
for {\bf compressible Rivlin}.
Finally, 
\be
 W = \underset{\smVOL}{\underbrace{\frac{K}{2}(J - 1)^2}} + 
 \underset{\smDEV}{\underbrace{C_{10} (\bar{I}_1 - 3)}},  
\ee
for {\bf compressible Neo-Hookean}.


